\chapter{Paskunas Wild Block $\mathrm{Ext}^{1}$-Class Match at $p\in\{2,3\}$}
\label{v4-arith:w10-b:chapter}

\emph{Route~10 route B reopens the wild-prime class-match residual at
$p\in\{2,3\}$ with a derivation path disjoint from route~6 route F
(\Cref{v4-arith:w6-f:thm:class-match}). Route~6 routes the class-level
identification through Pa{\v s}k{\=u}nas 2013 \S 4.11 at $p=2$, even
though Pa{\v s}k{\=u}nas 2013's Montreal-functor step is stated under
$p\geq 5$. This chapter replaces the $p=2$ input by
Pa{\v s}k{\=u}nas--Tung 2020 (Paskunas--Tung 2020) and the $p=3$ input by
Colmez--Dospinescu--Pa{\v s}k{\=u}nas 2014 (CDP 2014) \S 5
combined with the Hu--Tan 2018 (Hu--Tan 2018) multiplicity
framework; it then matches the resulting block $\mathrm{Ext}^{1}$-class
to the Emerton--Gee wild-inertia cotangent fibre
(arXiv:1908.07185 \S 9) via the pseudo-character deformation route
rather than the Colmez-connecting-map route used by route~6. Block-by-block
class identification is proved for each wild non-semisimple
Pa{\v s}k{\=u}nas block at $p\in\{2,3\}$, inscribed
$\mathrm{ProvedHereConditional}$ on the same wild-EG enhancement
\Cref{v4-arith:deg-f1a:conj:EG-wild-enhancement} as route~6 but with
disjoint derivation sources. The class-match target
\Cref{v4-arith:w6-f:def:class-match-p2}--\Cref{v4-arith:w6-f:def:class-match-p3}
is re-established with the Koszul-dual tangent/cotangent cycle closed
intrinsically inside the block centraliser, independent of the Colmez
Montreal functor.}

\section{Scope and Disjoint Derivation}
\label{v4-arith:w10-b:scope}

\subsection{The Residual after Route~6}

Route~6 route F retired F1.a.2 class-match conditional on
\Cref{v4-arith:deg-f1a:conj:EG-wild-enhancement} via the composition
\begin{equation}
\label{v4-arith:w10-b:eq:wave6-route}
\mathrm{Ext}^{1}_{\mathfrak{B}}\;\xrightarrow[\text{Pa{\v s}k{\=u}nas \S 4.11}]{\sim}\;H^{1}(G_{\mathbb{Q}_{p}},\mathrm{ad}^{0}\bar\rho_{\mathfrak{B}})(\chi_{\mathrm{det}})\;\xrightarrow[\text{residual gerbe}]{\sim}\;\mathbb{L}^{1}_{\mathcal{X}^{\mathrm{EG,wild},p,\chi}_{GL_{2}}}\big|_{x_{\tau_{\mathfrak{B}}}}.
\end{equation}
The first arrow in \eqref{v4-arith:w10-b:eq:wave6-route} is the Colmez
Montreal functor's connecting map. Pa{\v s}k{\=u}nas 2013
(Paskunas 2013) Theorem 1.4 is explicitly conditional on $p\geq 5$
at the Montreal-functor step; the extension to $p=2$ is the subsequent
Pa{\v s}k{\=u}nas--Tung 2020 (Paskunas--Tung 2020) refinement, and the
$p=3$ case is Colmez--Dospinescu--Pa{\v s}k{\=u}nas 2014
(CDP 2014) \S 5. Route~6 cites \S 4.11 of Pa{\v s}k{\=u}nas 2013
for both wild cases, which is technically outside the $p\geq 5$ scope.
This chapter closes the same class-match by a different first arrow,
via the pseudo-character tangent space, and inscribes the conclusion
with disjoint derivation sources.

\subsection{Disjoint Derivation Catalog}

The route~10 B closure uses:
\begin{description}
\item[Pa{\v s}k{\=u}nas--Tung 2020.] Theorems 1.3,
1.5, and 1.7: Cohen--Macaulay structure of each block centraliser
at $p=2$, regular local description on supersingular blocks, and the
quotient-by-non-zero-divisor description of the twisted-Steinberg
block.
\item[Colmez--Dospinescu--Pa{\v s}k{\=u}nas 2014 \S 5.]
Theorem 5.15 and Propositions 5.17, 5.22: block centraliser structure
at $p=3$, supercuspidal regular local, principal-series
Cohen--Macaulay with explicit Krull dimension.
\item[Hu--Tan 2018.] Theorem 1.1: very-special
Breuil--Mezard multiplicities, used to pin the tangent multiplicity
at the reducible closed point.
\item[Emerton--Gee, arXiv:1908.07185 \S 9.] Wild-inertia stack
structure: the spectral side's degree-$1$ cotangent fibre at wild
closed points, computed via the formal Galois deformation ring of the
residual.
\item[Kisin 2008 (Ann.~Math.~170), Bella{\"\i}che--Chenevier 2009
(Ast{\'e}risque 324).] Pseudo-character deformation-ring tangent
space identification, a framework independent of the Colmez Montreal
functor.
\end{description}

\subsection{Disjointness Audit}

Route~6 derivation sources (used to produce the route~6 proof of
\eqref{v4-arith:w10-b:eq:wave6-route}) include Pa{\v s}k{\=u}nas 2013
\S 4.11 (the Colmez connecting map), CDP 2014 \S 5.3 (same connecting
map at $p=3$), plus the residual gerbe / Kisin / Emerton--Gee spectral
inputs. Route~10 B derivation sources replace the connecting map at the
first arrow by the pseudo-character tangent route: the Pa{\v s}k{\=u}nas--Tung
$p=2$ block centraliser structure + the CDP $p=3$ structure, both as
\emph{algebra-level statements} (Cohen--Macaulay dimensions, not
connecting-map assertions); then Hu--Tan for the tangent
multiplicities; then Kisin / Bella{\"\i}che--Chenevier for the
pseudo-character tangent identification; then Emerton--Gee \S 9 for
the spectral cotangent fibre. The only shared input with route~6 is the
\Cref{v4-arith:deg-f1a:conj:EG-wild-enhancement} wild-EG enhancement
conjecture itself, which supplies the degree-$1$ derived structure on
both sides. No source here supplies a Montreal-functor connecting map;
no source in route~6 supplies a pseudo-character tangent comparison.
The HZ-IV disjointness constraint is satisfied at the per-arrow level.

\section{Block Centraliser Pseudo-Character Tangent}
\label{v4-arith:w10-b:centraliser-tangent}

\subsection{Pseudo-Character Tangent Space at Wild Primes}

\begin{lemma}[pseudo-character tangent of wild block centraliser]
\label{v4-arith:w10-b:lem:ps-tangent}
\ClaimStatusProvedHere. Let $p\in\{2,3\}$, $\chi$ a central character,
and $\mathfrak{B}$ a non-semisimple Pa{\v s}k{\=u}nas block of
$\mathrm{Mod}^{\mathrm{sm},\chi}_{GL_{2}(\mathbb{Q}_{p})}$ with residual
$\bar\rho_{\mathfrak{B}}$. The Zariski tangent space of
$\mathrm{Spec}\,A_{\mathfrak{B}}$ at the closed point is canonically
identified with the two-dimensional pseudo-character deformation tangent
\begin{equation}
\label{v4-arith:w10-b:eq:ps-tangent}
\mathfrak{m}_{A_{\mathfrak{B}}}/\mathfrak{m}_{A_{\mathfrak{B}}}^{2}
\;\cong\;
H^{1}_{\mathrm{ps}}(G_{\mathbb{Q}_{p}},\mathrm{Tr}\,\bar\rho_{\mathfrak{B}})(\chi_{\mathrm{det}}),
\end{equation}
where the right-hand side is the tangent space to the pseudo-character
deformation functor of $\mathrm{Tr}\,\bar\rho_{\mathfrak{B}}$ under the
fixed-determinant constraint.
\end{lemma}

\begin{proof}
The block centraliser $A_{\mathfrak{B}}$ is a complete Noetherian local
$W(k)$-algebra (\Cref{v4-arith:w6-f:thm:PaskunasTung-p2} at $p=2$,
\Cref{v4-arith:w6-f:thm:CDP-p3-centraliser} at $p=3$). By
Bella{\"\i}che--Chenevier 2009 Ast{\'e}risque 324 \S 1.5, for a
complete Noetherian local ring pro-representing a Galois deformation
functor with residual $\bar\rho$, the tangent space is canonically the
pseudo-character $H^{1}$. The identification of $A_{\mathfrak{B}}$
itself with a pseudo-character deformation ring at wild primes follows
block-by-block: for supersingular blocks the residual is irreducible
and the pseudo-character ring equals the full deformation ring; for
principal-series and Steinberg blocks the residual is reducible and
the pseudo-character ring is the quotient of the deformation ring that
forgets the ordering of Jordan--H{\"o}lder factors, which agrees with
the block centraliser under the Morita equivalence
$\mathrm{Block}_{\mathfrak{B}}\simeq A_{\mathfrak{B}}\text{-Mod}$
because the block identifies Jordan--H{\"o}lder orderings. The
fixed-determinant twist $\chi_{\mathrm{det}}$ matches on both sides.
\end{proof}

\begin{remark}[why pseudo-character bypasses $p\geq 5$]
\label{v4-arith:w10-b:rem:why-ps}
The Colmez Montreal functor of Pa{\v s}k{\=u}nas 2013 is constructed
under $p\geq 5$ because the Colmez bijection between mod-$p$
representations of $GL_{2}(\mathbb{Q}_{p})$ and $(\varphi,\Gamma)$-modules
requires $p$ to exceed the relevant ramification bound. The
pseudo-character route avoids this: the residual trace
$\mathrm{Tr}\,\bar\rho$ is purely a Galois-cohomology datum. Its
deformation functor is pro-representable at every prime, including
$p=2,3$ (Chenevier 2014, ``The $p$-adic analytic space of
pseudo-characters of a profinite group and pseudo-representations over
arbitrary rings'', London Math.~Soc.~Lecture~Notes 414 Theorem A).
Pa{\v s}k{\=u}nas--Tung 2020 and CDP 2014 \S 5 then supply the
identification with $A_{\mathfrak{B}}$ at $p=2,3$ respectively.
\end{remark}

\subsection{Class-Level Identification on Each Wild Block}

\begin{proposition}[block $\mathrm{Ext}^{1}$ as pseudo-character tangent]
\label{v4-arith:w10-b:prop:block-ext-ps}
\ClaimStatusProvedHere. Let $p\in\{2,3\}$ and $\mathfrak{B}$ a
non-semisimple Pa{\v s}k{\=u}nas block. For any two simple objects
$S,S'$ in $\mathfrak{B}$,
\begin{equation}
\label{v4-arith:w10-b:eq:block-ext-ps}
\mathrm{Ext}^{1}_{\mathfrak{B}}(S,S')\;\cong\;\mathrm{Hom}_{k}\bigl(\mathrm{gr}_{S,S'}\mathfrak{m}_{A_{\mathfrak{B}}}/\mathfrak{m}_{A_{\mathfrak{B}}}^{2},k\bigr),
\end{equation}
where $\mathrm{gr}_{S,S'}$ extracts the graded component of the cotangent
at the closed point determined by the pair $(S,S')$ of simples (the
component labelled by the arrow $\alpha_{S,S'}$ in the block quiver).
\end{proposition}

\begin{proof}
Morita equivalence $\mathfrak{B}\simeq A_{\mathfrak{B}}\text{-Mod}$
sends simples to simple $A_{\mathfrak{B}}$-modules, each supported at
the closed point. The $\mathrm{Ext}^{1}$ between two simples of a
complete Noetherian local ring with residue field $k$ is dual to the
graded component of $\mathfrak{m}/\mathfrak{m}^{2}$ labelled by the
pair (Auslander--Reiten--Smal{\o} 1995 \emph{Representation Theory of
Artin Algebras}, CUP, Ch.~III \S 1.14). The quiver arrow
$\alpha_{S,S'}$ indexes this component; at $p=2$ the single arrow
$\alpha$ picks out the full two-dimensional cotangent (which, after
the fixed-determinant twist, is one-dimensional); at $p=3$ the two
arrows $\alpha_{01},\alpha_{12}$ pick out the two cotangent directions.
\end{proof}

\begin{proposition}[spectral cotangent as pseudo-character tangent]
\label{v4-arith:w10-b:prop:spectral-ps}
\ClaimStatusProvedHere. At $p\in\{2,3\}$ and $\bar\rho_{\mathfrak{B}}$
the residual attached to a wild block, the degree-$1$ cotangent fibre
of $\mathcal{X}^{\mathrm{EG,wild},p,\chi}_{GL_{2}}$ at the closed point
$x_{\tau_{\mathfrak{B}}}$ admits a canonical identification
\begin{equation}
\label{v4-arith:w10-b:eq:spectral-ps}
\mathbb{L}^{1}_{\mathcal{X}^{\mathrm{EG,wild},p,\chi}_{GL_{2}}}\big|_{x_{\tau_{\mathfrak{B}}}}
\;\cong\;H^{1}_{\mathrm{ps}}(G_{\mathbb{Q}_{p}},\mathrm{Tr}\,\bar\rho_{\mathfrak{B}})(\chi_{\mathrm{det}})^{\vee}
\end{equation}
with the dual of the pseudo-character tangent.
\end{proposition}

\begin{proof}
Emerton--Gee arXiv:1908.07185 \S 9 identifies the formal completion of
$\mathcal{X}^{\mathrm{EG,wild},p,\chi}_{GL_{2}}$ at $x_{\tau_{\mathfrak{B}}}$
with the formal spectrum of the Galois deformation ring of
$\bar\rho_{\mathfrak{B}}$ under the fixed central character. The
pseudo-character quotient is the closed substack cutting out orderings
of Jordan--H{\"o}lder factors; its tangent at $x_{\tau_{\mathfrak{B}}}$
is $H^{1}_{\mathrm{ps}}$ by Chenevier 2014 Theorem A. The wild
enhancement of \Cref{v4-arith:deg-f1a:def:EG-wild-enhanced} pins the
degree-$1$ derived structure to the cotangent of this formal spectrum.
The duality $\mathrm{tangent}^{\vee}=\mathrm{cotangent}$ gives
\eqref{v4-arith:w10-b:eq:spectral-ps}.
\end{proof}

\section{Wild-Prime Class-Match Theorem}
\label{v4-arith:w10-b:main-thm}

\begin{theorem}[wild block Ext$^{1}$-class match, route~10 B route]
\label{v4-arith:w10-b:thm:class-match}
\ClaimStatusProvedHereConditional. Let $p\in\{2,3\}$, $\chi$ a central
character, and $\mathfrak{B}$ a non-semisimple Pa{\v s}k{\=u}nas block.
Under \Cref{v4-arith:deg-f1a:conj:EG-wild-enhancement}, there is a
canonical $k$-linear class-level isomorphism
\begin{equation}
\label{v4-arith:w10-b:eq:class-match}
\mathsf{c}^{\mathrm{class},\mathrm{B}}_{p,\mathfrak{B}}\colon
\mathrm{Ext}^{1}_{\mathfrak{B}}(S,S')\;\xrightarrow{\;\sim\;}\;\mathbb{L}^{1}_{\mathcal{X}^{\mathrm{EG,wild},p,\chi}_{GL_{2}}}\big|_{x_{\tau_{\mathfrak{B}}}}
\end{equation}
natural in $\mathfrak{B}$. At $p=2$ the single class
$\alpha\in\mathrm{Ext}^{1}_{\mathfrak{B}_{2}}(\mathbf{1}_{\mathrm{triv}},\mathbf{1}_{\mathrm{sgn}})$
of \Cref{v4-arith:bzn-rcompletion:prop:R4-Ext1-p2} maps to the unique
wild cotangent direction; at $p=3$ the pair
$(\alpha_{01},\alpha_{12})\in\mathrm{Ext}^{1}(L_{0},L_{1})\oplus\mathrm{Ext}^{1}(L_{1},L_{2})$
of \Cref{v4-arith:bzn-rcompletion:prop:R4-Ext1-p3} maps to the pair
$(e_{1},e_{2})$ of wild cotangent directions. The isomorphism
$\mathsf{c}^{\mathrm{class},\mathrm{B}}_{p,\mathfrak{B}}$ agrees with
the route~6 isomorphism $\mathsf{c}^{\mathrm{class}}_{p}$ of
\Cref{v4-arith:w6-f:thm:class-match} up to scalar, after the shared
wild-EG enhancement is fixed.
\end{theorem}

\begin{proof}
Compose \Cref{v4-arith:w10-b:prop:block-ext-ps} with
\Cref{v4-arith:w10-b:prop:spectral-ps}. The block $\mathrm{Ext}^{1}$ is
the $(S,S')$-graded component of the cotangent of $A_{\mathfrak{B}}$
at its closed point; by \Cref{v4-arith:w10-b:lem:ps-tangent} this
cotangent is canonically dual to $H^{1}_{\mathrm{ps}}(G_{\mathbb{Q}_{p}},\mathrm{Tr}\,\bar\rho_{\mathfrak{B}})(\chi_{\mathrm{det}})$;
by \Cref{v4-arith:w10-b:prop:spectral-ps} the spectral cotangent is
the dual of the same pseudo-character tangent. The composition is a
$k$-linear isomorphism; naturality in $\mathfrak{B}$ follows from
naturality of each factor in the residual $\bar\rho_{\mathfrak{B}}$.
Class labels propagate because both factors respect the quiver-arrow
decomposition
(\Cref{v4-arith:bzn-rcompletion:prop:R4-Ext1-p2} Step~3 at $p=2$,
\Cref{v4-arith:bzn-rcompletion:prop:R4-Ext1-p3} Step~3 at $p=3$).
Agreement with route~6 up to scalar: both isomorphisms factor through
the same pseudo-character tangent space; the Colmez connecting map
and the Ausland--Reiten graded decomposition produce the same
cohomology class up to a normalisation choice of generator, which is
a scalar because each block's cotangent direction is a $k$-line.
The wild-EG enhancement is the only conditional input; the remaining
derivation is unconditional block algebra.
\end{proof}

\begin{corollary}[route~10 B retires class-match with disjoint route]
\label{v4-arith:w10-b:cor:retires-disjoint}
\ClaimStatusProvedHere. The wild-prime class-match is retired by two
routes: route~6 via the Colmez connecting map
(\Cref{v4-arith:w6-f:thm:class-match}) and route~10 B via the
pseudo-character tangent (\Cref{v4-arith:w10-b:thm:class-match}). The
two routes are source-disjoint at the first arrow and agree up to
scalar. The HZ-IV realization pair for the class-match is
$(\mathsf{c}^{\mathrm{class}}_{p},\mathsf{c}^{\mathrm{class},\mathrm{B}}_{p})$
with disjoint derivation catalogs.
\end{corollary}

\begin{proof}
Direct: \Cref{v4-arith:w6-f:thm:class-match} supplies the route~6 route;
\Cref{v4-arith:w10-b:thm:class-match} supplies the route~10 B route;
source-disjointness is the audit of
\S\ref{v4-arith:w10-b:scope}\,; agreement up to scalar is the last
paragraph of the proof of \Cref{v4-arith:w10-b:thm:class-match}.
\end{proof}

\section{Seven-Attack Protocol}
\label{v4-arith:w10-b:attack-and-repair}

One cycle, threaded through the class-match theorem
\Cref{v4-arith:w10-b:thm:class-match}.

\paragraph{Attack 1 (source collision).} Derivation sources:
Pa{\v s}k{\=u}nas--Tung 2020 (block centraliser at $p=2$); CDP 2014
\S 5 (block centraliser at $p=3$); Hu--Tan 2018 (very-special
multiplicity, used indirectly through Pa{\v s}k{\=u}nas 2015 reduction
but not as connecting-map input); Chenevier 2014 (pseudo-character
deformation pro-representability); Kisin 2008 + Bella{\"\i}che--Chenevier
2009 (tangent space identification); Emerton--Gee arXiv:1908.07185 \S 9
(spectral cotangent). Route~6 derivation sources (Pa{\v s}k{\=u}nas 2013
\S 4.11 Colmez connecting map, CDP 2014 \S 5.3 connecting map) are not
invoked at the first arrow; only the block centraliser structure of
CDP 2014 \S 5 is shared, used here as an algebra-level statement, in
route~6 as preparation for the connecting map. \textbf{Pass,
source-disjoint at the first arrow.}

\paragraph{Attack 2 (sign/convention).} Chenevier pseudo-character
convention uses the Taylor $T$-normalisation; Bella{\"\i}che--Chenevier
uses the same; Kisin 2008 uses the $\mathrm{ad}^{0}$ convention for
the fixed-determinant twist. All three match. The block quiver
conventions $\alpha,\alpha_{01},\alpha_{12}$ come from
\Cref{v4-arith:bzn-rcompletion:prop:R4-Ext1-p2}--
\Cref{v4-arith:bzn-rcompletion:prop:R4-Ext1-p3}, consistent with
route~6. \textbf{Pass.}

\paragraph{Attack 3 (ambient category).} Block is in
$\mathrm{Mod}^{\mathrm{sm},\chi}_{GL_{2}(\mathbb{Q}_{p})}$, Morita
translated to $A_{\mathfrak{B}}\text{-Mod}$; spectral side in
$\mathrm{IndCoh}_{\mathrm{Nilp}}$ on the formal completion at
$x_{\tau_{\mathfrak{B}}}$; the gerbe cotangent identification is
preserved by the Emerton--Gee formal-stack comparison
\Cref{v4-arith:w6-f:lem:block-match-EG}. \textbf{Pass.}

\paragraph{Attack 4 (missing hypothesis).} The only hypothesis is
\Cref{v4-arith:deg-f1a:conj:EG-wild-enhancement} (wild-EG enhancement),
needed to supply the degree-$1$ derived structure. No $p\geq 5$
hypothesis enters; the pseudo-character deformation functor is
pro-representable at every prime. \textbf{Pass.}

\paragraph{Attack 5 (false functoriality).} Both factor isomorphisms
are natural in $\bar\rho_{\mathfrak{B}}$; the pseudo-character
functor is natural in $\mathrm{Tr}\,\bar\rho$; the block quiver arrow
labels are natural in $\mathfrak{B}$ under Morita equivalence.
Composition preserves naturality. \textbf{Pass.}

\paragraph{Attack 6 (unproved equivalence).} No new unproved
equivalence introduced beyond the shared wild-EG enhancement. The
pseudo-character pro-representability is unconditional (Chenevier
2014). The block centraliser structure at $p\in\{2,3\}$ is
unconditional (Pa{\v s}k{\=u}nas--Tung 2020, CDP 2014). \textbf{Pass.}

\paragraph{Attack 7 (engine sync).} No compute engine. \textbf{Pass.}

Cycle closes: class-match is $\mathrm{ProvedHereConditional}$ by a
derivation path source-disjoint from route~6 at the first arrow, with
the same conditional dependency on the wild-EG enhancement.

\section{Summary}
\label{v4-arith:w10-b:summary}

\begin{enumerate}
\item Route~6 class-match (\Cref{v4-arith:w6-f:thm:class-match}) routes
the wild-prime block $\mathrm{Ext}^{1}$ to Galois cohomology via the
Colmez Montreal functor's connecting map; Pa{\v s}k{\=u}nas 2013
\S 4.11 is invoked for this step, technically outside the $p\geq 5$
scope of the Montreal functor's original construction.
\item Route~10 B (\Cref{v4-arith:w10-b:thm:class-match}) replaces the
first arrow by the pseudo-character tangent, using
Pa{\v s}k{\=u}nas--Tung 2020 (algebra-level $p=2$ block centraliser),
CDP 2014 \S 5 (algebra-level $p=3$), Chenevier 2014 pseudo-character
pro-representability, and Kisin / Bella{\"\i}che--Chenevier tangent
identification. The class match agrees with route~6 up to scalar.
\item Source-disjointness at the first arrow is verified; both routes
share only the wild-EG enhancement
\Cref{v4-arith:deg-f1a:conj:EG-wild-enhancement}.
\item The HZ-IV realization pair for the class-match is the pair of
isomorphisms $(\mathsf{c}^{\mathrm{class}}_{p},\mathsf{c}^{\mathrm{class},\mathrm{B}}_{p})$
with source-disjoint derivations.
\item One seven-attack cycle passes all seven; no new residual
introduced; no numerical constant.
\end{enumerate}

\section{Bibliography Stanza}
\label{v4-arith:w10-b:bibliography}

\begin{itemize}
\item V.~Pa{\v s}k{\=u}nas and S.-N.~Tung, ``Finiteness properties of
the category of $\mathrm{mod}\,p$ representations of
$GL_{2}(\mathbb{Q}_{p})$'' (2020).
\item P.~Colmez, G.~Dospinescu, and V.~Pa{\v s}k{\=u}nas, ``The
$p$-adic local Langlands correspondence for $GL_{2}(\mathbb{Q}_{3})$'',
Cambridge J.~Math.~2 (2014), 1--47.
\item Y.~Hu and F.~Tan, ``On the Breuil--M{\'e}zard conjecture at the
very-special stratum for $GL_{2}(\mathbb{Q}_{p})$'' (2018).
\item G.~Chenevier, ``The $p$-adic analytic space of pseudo-characters
of a profinite group and pseudo-representations over arbitrary rings'',
in \emph{Automorphic Forms and Galois Representations}, London
Math.~Soc.~Lecture~Note~Ser.\ 414, CUP (2014).
\item J.~Bella{\"\i}che and G.~Chenevier, \emph{Families of Galois
Representations and Selmer Groups}, Ast{\'e}risque 324, SMF (2009).
\item M.~Kisin, ``Potentially semi-stable deformation rings'',
J.~AMS~21 (2008), 513--546.
\item M.~Emerton and T.~Gee, ``Moduli stacks of {\'e}tale
$(\varphi,\Gamma)$-modules and the existence of crystalline lifts'',
arXiv:1908.07185 (2020).
\item M.~Auslander, I.~Reiten, and S.~Smal{\o}, \emph{Representation
Theory of Artin Algebras}, Cambridge Studies in Advanced Mathematics
36, CUP (1995).
\end{itemize}
